\documentclass[12pt]{article}
\usepackage[margin=1in]{geometry}
\usepackage{graphicx}
\usepackage{booktabs}
\usepackage{amsmath}
\title{Replicating <one specific result> from <Author (Year)>}
\author{Your Name}
\date{\today}
\begin{document}
\maketitle
\section{The Original Paper}
% 1-2 paragraphs:
% - Full citation
% - The paper's central research question
% - The empirical strategy in one or two sentences
\section{The Result of Interest}
% 1 paragraph:
% - The specific table cell / figure panel / number you are replicating
% - Why this particular result is the one worth reproducing
% (Is it the headline number? The cleanest identification? Most cited?)
\section{Data and Methods}
% 1-2 paragraphs:
% - Where the data came from
% - What preprocess.py does
% - What analysis.py does
\section{Replication Results}
% Include the table or figure produced by your code:
\begin{table}[h]
\centering
\caption{Replication of Summary Employment Statistics (Table 5),
         from Jardim et al.\ (2022)}
\begin{tabular}{lccc}
\hline
 & 2014 Q2 & 2016 Q2 & Change \\
\hline
Jobs paying $<\$19$/hr    & $90,757$ & $89,188$ & $-1.7\%$ \\
Hours worked $<\$19$/hr   & $36,450,720$ & $35,466,512$ & $-2.7\%$ \\
Avg.\ wage $<\$19$/hr     & $\$14.19$ & $\$15.00$ & $5.7\%$ \\
\hline
\end{tabular}
\label{tab:main}
\end{table}
% and / or
\begin{figure}[h]
 \centering
 \includegraphics[width=0.85\textwidth]
{../output/figures/main_figure.png}
 \caption{...}
\end{figure}
% 1-2 paragraphs:
% - Compare your numbers to the paper's
% - If they match: state that clearly and note any minor discrepancies
% - If they don't match: be specific about the gap and your best guess at
why
% (data version, sample restrictions, software differences, etc.)
\section{What I Learned}
% 1-2 paragraphs (this section is REQUIRED and graded):
% - What was easiest about the replication?
% - What was hardest? (data prep is almost always the answer)
% - What did you learn about reproducibility / about this paper that you
% could only learn by doing the replication?
replication_final.md 2026-05-21
6 / 9
% - If you were the original author, what would have made your work easier
% to replicate?
\begin{thebibliography}{9}
\bibitem{paper} <full citation of the paper you replicated>
\end{thebibliography}
\end{document}
